\section{Preliminaries on permutation actions on a group}
Let $G$ be a finite group.
For each $g\in G$, let $\lambda_g\in \Sym{G}$ 
be left multiplication with $g$
(so $\lambda_g(x)= gx$),
and $\rho_g$ be right multiplication
with $g^{-1}$, that is, $\rho_g(x)= xg^{-1}$.
Thus $g\mapsto \lambda_g$ and $g\mapsto \rho_g$ are 
the left and right regular permutation action.
Also, let $\iota\in \Sym{G}$ be the map that 
inverts elements (so $\iota(x)= x^{-1}$ for all $x\in G$).
Let $\Gamma(G)\leq \Sym{G}$ be the group generated by all these
elements:
\[ \Gamma(G):= \erz{\lambda_g, \rho_g, \iota
                   \mid g\in G}.
\]
To describe $\Gamma(G)$, we need the
\emph{wreath product} $G\wr C_2$ 
of $G$ with a cyclic group~$C_2=\erz{s}$
of order $2$.
Recall that this is the semidirect product of
$G\times G$ with $C_2$, where $s$ acts on 
$G\times G$ by exchanging coordinates:
$(g,h)^s = (h,g)$ for $g$, $h\in G$.
Then:

\begin{lemma}
  If $G$ is not an elementary abelian $2$-group,
  then $\Gamma(G) \iso (G\wr C_2)/Z$,
  where $Z = \{(z,z)\in G \times G \mid z \in \Zentr(G)\}$.
\end{lemma}

\begin{proof}
  We have that $\lambda(G)$ and $\rho(G)$ centralize each other,
  and $(\lambda_g)^{\iota} = \rho_g$.
  Thus sending $(g,h)\in G\times G$ to
  $\lambda_g \rho_h$ and $s\in C_2 = \{1,s\}$
  to $\iota$ defines a surjective group homomorphism
  $G\wr C_2 \to \Gamma(G)$
  with $Z$ in the kernel.
  
  Suppose $\lambda_g \rho_h = \Id_G$.
  Then $gxh^{-1} = x$ for all $x\in G$.
  Taking $x=1$ yields $g=h$, 
  and it follows that  $g\in \Zentr(G)$.
  
  Now assume $\lambda_g \rho_h \iota = \Id$.
  Then $gx^{-1}h^{-1} = x$ for all $x\in G$,
  and $x=1$ yields $g=h$.
  Moreover, we have
  $xy = g (xy)^{-1} g^{-1}
      = g y^{-1} g^{-1} \, g x^{-1} g^{-1}
      = y x$
  for all $x$, $y\in G$.
  Thus $G$ must be abelian in this case, and
  $x^{-1} = x$ for all $x\in G$.
    
  So when $G$ is not an elementary abelian $2$-group,
  such an element can not be in the kernel of the action
  of $G\wr C_2$ on $G$.
  This shows the result.
\end{proof}

In the proof of Theorem~\ref{main:birkhoff_unique}, 
we need the fact
that $\Gamma(G)$ contains no pair of commuting,
regular subgroups other than $\lambda(G)$
and $\rho(G)$, when $G = S_n$ and $n\geq 4$.
The exception in Theorem~\ref{main:birkhoff_unique}
for $n=3$
comes from the fact that 
in $\Gamma(S_3)$, we have other pairs of commuting, regular
subgroups, namely $U=V = C_2 \times C_3$
and $U=V= C_3\times C_2$.
Notice that we do not assume that the commuting, regular subgroups
$U$, $V$
of $\Gamma(G)$ have trivial intersection.
If one assumes $U\cap V=1$, one can give a somewhat 
shorter proof that 
$\{U,V\} = \{ \lambda(G), \rho(G)\}$
for almost simple groups $G$,
but we need the stronger statement for the proof of 
Theorem~\ref{main:birkhoff_unique}.

The most elegant and elementary way to prove 
that $\lambda(G)$ and $\rho(G)$ form the only pair
of commuting regular subgroups of $\Gamma(G)$
(when $G= S_n$, $n\geq 4$),
seems to be to use a general argument due to
Chermak and Delgado~\cite{ChermakDelgado89}.
Let $G$ be an arbitrary finite group.
Following Isaacs~\cite[\S~1G]{IsaacsFGT},
we call $m_G(H) := \card{H}\card{\Centrl_G(H)}$
the \emph{Chermak-Delgado measure} of the subgroup
$H\leq G$.

\begin{lemma}\cite[Theorem~1.44]{IsaacsFGT}
  \label{l:cdlatt}
  Let $G$ be a finite  group and let
  $\mathcal{L} = \mathcal{L}(G)$ be the set of subgroups
  for which the Chermak-Delgado measure is as large as 
  possible.
  Then for $H$, $K\in \mathcal{L}$, we have
  $H\cap K\in \mathcal{L}$,
  $\erz{H,K} =HK=KH \in \mathcal{L}$, 
  and $\Centrl_G(H)\in \mathcal{L}$.
\end{lemma}

The \emph{Chermak-Delgado lattice} of $G$ is by definition
the set of all subgroups of $G$ for which 
the Chermak-Delgado measure is maximized.
The last result tells us that this is indeed a sublattice
of the lattice of all subgroups of $G$.
We need the following, which is probably well known:

\begin{cor}
  Any member of the Chermak-Delgado lattice
  of a finite group $G$ is subnormal in $G$.
\end{cor}

\begin{proof}
If $H$ is a member of the Chermak-Delgado lattice of $G$,
then any conjugate $H^g$ is also in the 
Chermak-Delgado lattice,
and so $HH^g = H^g H$ by Lemma~\ref{l:cdlatt}.
But subgroups $H\leq G$ with $HH^g = H^gH$ for all $g\in G$
are subnormal~\cite[Theorem~2.8]{IsaacsFGT}.
\end{proof}

\begin{lemma}\label{l:sn_cent_est}
    Suppose that $G$ is almost simple
    (that is, $G$ has a nonabelian simple
    socle). 
    Then $\card{U}\card{\Centrl_{G}(U)} \leq \card{G}$
    for any subgroup $U\leq G$,
    and equality holds if and only if
    $U = \{1\}$ or $U=G$.
    In particular, this holds for 
    $G= S_n$, $n\geq 5$.
    The conclusion is also true for $G=S_4$.
\end{lemma}

\begin{proof}
  Suppose that $1\neq H$ is a member of
  the Chermak-Delgado lattice.
  Then $H$ is subnormal and thus contains the nonabelian
  simple socle of $G$.
  It follows that $\Zentr(H)=1 = H \cap \Centrl_G(H)$.
  Since $\Centrl_G(H)$ is also a member of the Chermak-Delgado lattice,
  we must have $\Centrl_G(H)=1$.
  Since $\card{H}\card{\Centrl_G(H)} = \card{H} \leq \card{G}$
  was supposed to be maximal possible, we see that
  $H=G$.
  Thus the Chermak-Delgado lattice contains exactly the groups
  $1$ and $G$ itself,
  and the first assertion follows.
  The case $G=S_4$ is a simple verification.
\end{proof}

We will need the following application (for $G=S_n$):

\begin{lemma}\label{l:snwrc2_reg}
  Let $G$ be a group such that the Chermak-Delgado lattice of $G$
  contains exactly the groups $1$ and $G$.
  Then $\lambda(G)$, $\rho(G)$ is the only pair of
  commuting, regular subgroups of\/ $\Gamma(G)$.
\end{lemma}

\begin{proof}
    Notice that $\Zentr(G)=\{1\}$, since otherwise
    $m_G(\Zentr(G))=\card{\Zentr(G)}\card{G} > \card{G} = m_G(1)$.
    Thus $\Gamma(G)\iso G\wr C_2$
    and $\lambda(G)\rho(G) \iso G\times G$.
    
    We first show that a regular subgroup $U$ of 
    $\Gamma(G)$ is contained in the normal subgroup
    $\lambda(G) \rho(G)$.
    Otherwise, $U$
    contains an element $u = \lambda_g \rho_h \iota$
    sending $x\in G$ to $ g x^{-1} h^{-1}$.
    Then $u^2$ sends $x$ to $gh x g^{-1}h^{-1}$,
    and in particular fixes $g$.
    By regularity, we must have $u^2 = \Id_G$. 
    This implies 
    $gh = hg $ and 
    $gh\in \Zentr(G) = \{1\}$.
    Thus $u$ sends $x$ to $g x^{-1}g$,
    and so fixes $g$, too, 
    which contradicts the regularity. 
    This shows that $U \leq \lambda(G) \rho(G)$.
    
    Since $\lambda(G) \rho(G)\iso G\times G$, we may work
    in $G\times G$ from now on.
    Suppose that $U$ and $V\leq G\times G$ 
    both have size $\card{G}$,
    and commute with each other.
    Let $U_L$ be the projection of $U$ 
    onto the first component, 
    that is, the subgroup of elements
    $g\in G$ such that there is an $h\in G$
    with $(g,h)\in U$.
    Let $U_R$ be the projection of $U$ on the second component.
    With this notation,
    $\Centrl_{G\times G}(U) = \Centrl_{G}(U_L)\times \Centrl_{G}(U_R)$.
    Thus 
    \[ \card{G}^2 
       = \card{U}\card{V}
       \leq \card{U_L}\card{U_R}
            \card{\Centrl_{G}(U_L)}\card{\Centrl_{G}(U_R)}
       \leq \card{G}^2,
    \]
    where the last inequality follows from our assumption
    on the Chermak-Delgado lattice of $G$.
    Thus equality holds, and it follows also that
    $U_L$ and $U_R$ are trivial or the group $G$ itself.
    Since both $U$ and $V$ have size $\card{G}$,
    it follows that
    $\{U,V\} = \{G\times 1, 1\times G\}$.
\end{proof}

\begin{cor}\label{c:normlgamma}
  Let $G$ be a group such that the Chermak-Delgado lattice of $G$
  contains exactly the groups $1$ and $G$.
  Then $\Norml_{\Sym{G}}(\Gamma(G)) = (\Aut G)\Gamma(G)$.    
\end{cor}

\begin{proof}
  Let $\pi \in \Norml_{\Sym{G}}(\Gamma(G))$.
  Then $\lambda(G)^{\pi}$ and $\rho(G)^{\pi}$
  are commuting regular subgroups of $\Gamma(G)$,
  and thus 
  $\{\lambda(G)^{\pi},\rho(G)^{\pi}\}
   = \{ \lambda(G), \rho(G)\} $.
  Since $\lambda(G)$ and $\rho(G)$ are conjugate in 
  $\Gamma(G)$, 
  we may assume that $\lambda(G)^{\pi}= \lambda(G)$.
  Thus $\pi \lambda_g \pi^{-1} = \lambda_{\alpha g}$
  for some bijection $\alpha\colon G\to G$.
  Clearly, $\alpha$ is a group automorphism.
  
  As $\lambda(G)$ acts transitively on $G$, we may assume
  $\pi(1) = 1$. 
  But then 
  $\pi(g) = \pi\lambda_g \pi^{-1}(1)
          = \lambda_{\alpha g}(1)
          = \alpha(g)$,
  so $\pi\in \Aut G$.
\end{proof}

The conclusion of this corollary is also true for some other
groups (for example, $G=S_3$),
but not for all groups (for example, $G= S_3 \times S_3$).

